\documentclass[french]{book}
\usepackage{teach}
\usepackage{verbatim}
\usepackage{amsmath,amsfonts,amssymb}
\usepackage{pgfplots}
\usepackage{mwe}
\RequirePackage{graphicx}
\RequirePackage{ifpdf}
 \ifpdf
   \DeclareGraphicsRule{*}{mps}{*}{}
 \fi


  \usepackage{fancyhdr}

  \usepackage{amsmath}
  \usepackage{amssymb}
  \usepackage{amsfonts}
  \usepackage{amsthm}
  \usepackage{mathrsfs}

  \usepackage{array}
  \usepackage{multicol}
  \setlength{\multicolsep}{3pt}

  \usepackage{graphicx}
  \graphicspath{{Figures/}}
  \usepackage{pstricks,pst-plot,pst-text,pst-tree,pst-eucl}
  \usepackage[all]{xy}

  \usepackage{fancybox}
  \usepackage{color}

%  \usepackage[frenchb]{babel}
  \usepackage{hyperref}

\usepackage{subfig}
\pgfplotsset{compat=newest}
\usepackage[all]{xy}
%\usepackage{geometry}
%\geometry{top=1cm, bottom=1cm, left=2cm, right=2cm}
\usepackage[utf8]{inputenc}
\usepackage[french]{babel}
\redefinecolor{thm}{title}{bgcolor}{alizarin}
\redefinecolor{prop}{title}{bgcolor}{meatbrown}
\redefinecolor{defn}{title}{bgcolor}{olivine}
\definecolor{alizarin}{rgb}{0.82, 0.1, 0.26}
\definecolor{meatbrown}{rgb}{0.9, 0.72, 0.23}
\definecolor{olivine}{rgb}{0.6, 0.73, 0.45}
\usepackage{mathrsfs}
%%
\newcommand{\R}{\mathbb {R}}
%\newcommand{\C}{\mathds {C}}
\newcommand{\N}{\mathbb {N}}
\newcommand{\Z}{\mathbb {Z}}
\newcommand{\Q}{\mathbb {Q}}
\newcommand{\D}{\mathbb {D}}
%%
\newcommand{\se}{\geq}
\newcommand{\ie}{\leq}
\newcommand{\tm}{\times}
%\newcommand{\ell}{l}
%\newcommand{\ell'}{l'}
%%
\newcommand{\ds}{\displaystyle}
\newcommand{\limn}{\lim_{n\rightarrow +\infty}}
\newcommand{\limo}[1][x]{\ds\lim_{#1\rightarrow 0}}
\newcommand{\limite}[2][x]{\ds\lim_{#1\rightarrow #2}}
\newcommand{\limpinf}[1][x]{\ds\lim_{#1\rightarrow +\infty}}
\newcommand{\limminf}[1][x]{\ds\lim_{#1\rightarrow -\infty}}
\newcommand{\limiteg}[2][x]{\ds\lim_{#1\xrightarrow{<}#2}}
\newcommand{\limited}[2][x]{\ds\lim_{#1\xrightarrow{>}#2}}
%%
\newcommand{\vect}[1]{\overrightarrow{#1}}
\newcommand{\oij}{(\textit{O},{\overrightarrow{i}},{\overrightarrow{j}})}
%
\newcommand{\cp}[2]{%
        \begin{pmatrix}
        #1\\
        #2
        \end{pmatrix}%
        }
%
\newcommand{\calb}{\mathscr{B}}
\newcommand{\calc}{\mathscr{C}}
\newcommand{\cald}{\mathscr{D}}
\newcommand{\cale}{\mathscr{E}}
\newcommand{\calf}{\mathscr{F}}
\newcommand{\calp}{\mathscr{P}}
\newcommand{\cals}{\mathscr{S}}
\newcommand{\calt}{\mathscr{T}}
%
\newcommand{\suite}[1][u]{\left( #1_{n} \right)_{n\in\N}}
\newcommand{\suitep}[1][u]{\left( #1_{n} \right)_{n\in\N^{*}}}
\usetikzlibrary{arrows}

\newcommand{\poly}[1][a]{#1_0+#1_1x+\cdots+#1_nx^n}


\begin{document}
\tableofcontents
\
\chapter{Continuité}
\section{Fonction continue en un point, sur un intervalle}

\begin{defn}
 Soit $f$ une fonction définie sur un intervalle  $I$ et $a$ un nombre réel de $I$.
\begin{itemize}
 \item[$\bullet$] La fonction $f$ est dite \emph{continue en $a$} si elle admet en $a$ une limite égale à $f(a)$, c'est  à dire si $\lim_{x\rightarrow a}f(x)=f(a)$.
\item[$\bullet$] $f$ est dite \emph{continue sur $I$} si elle est continue en tout réel $a\in I$.
\end{itemize}
\end{defn}

\begin{exemple}
 Ci-dessous la fonction $f$, définie par $f(x)=(x+2)(x-1)(x+3)$ sur $\mathbb{R}$, est continue sur $]-\infty;+\infty[$ alors que la fonction $g$ définie sur $]-\infty;-1]$ par $g(x)=f(x)$ et sur $]-1;+\infty[$ par $g(x)=f(x)-1$ n'est pas continue en $-1$.
%% Graphique manquant
\end{exemple}


\begin{prop}
 Si $f$ est une fonction définie sur un intervalle $I$ et dérivable en un réel $a\in I$, alors $f$ est continue en $a$. En particulier, si $f$ est dérivable sur $I$, alors $f$ est continue sur $I$.
\end{prop}

\begin{rem}
 La réciproque n'est pas vraie comme le montre le contre exemple de la fonction racine carré en 0.
\end{rem}


\begin{prop}[continuité des fonctions de référence]
\begin{itemize}
 \item[$\bullet$] $x\mapsto \exp(x)$, les fonctions polynômes et la fonction valeur absolue sont continues sur $]-\infty;+\infty[$ ;
\item[$\bullet$] $x\mapsto \sqrt{x}$ est continue sur $[0;+\infty[$ (alors qu'elle n'est dérivable que sur $]0;+\infty[)$ ;
\item[$\bullet$] la fonction inverse est continue sur $]-\infty;0[$ et sur $]0;+\infty[$.
\end{itemize}
\end{prop}

\begin{demo}
 \'A l'exception du cas de la fonction racine carré et de la valeur absolue, la dérivabilité de chacune des fonctions suffit pour assurer la continuité.\\
Pour ce qui concerne la fonction racine carré, elle est dérivable donc continue sur $]0;+\infty[$. En outre $\sqrt{0}=0$ et pour tout $x$ réel tel que $0<x<1$, on sait que $0\leq\sqrt{x}\leq x$. Donc de $\lim_{x\rightarrow 0}x=0$ on déduit $\lim_{x\rightarrow 0}\sqrt{x}=0$ d'où la continuité en 0. Démonstration analogue pour la fonction valeur absolue.
\end{demo}

\section{Fonction dérivable $ \Longrightarrow $ fonction continue }

\begin{prop}
Soit $f$ une fonction dérivable en $a$ alors $f$ est continue en $a$.
\end{prop}

\begin{demo}
Soit $f$ une fonction dérivable en $a$, alors le nombre $f'(a)$ existe, par suite $\limite{a} (x-a)f'(a)= 0$  nécéssairement.

D'autre part $f'(a)= \limite{a} \dfrac{f(x)-f(a)}{x-a}$

D'où
 $$\limite{a} (x-a)f'(a) = \limite{a} (x-a)  \limite{a} \dfrac{f(x)-f(a)}{x-a} =  \limite{a} (x-a) \dfrac{f(x)-f(a)}{x-a} = \limite{a} f(x)-f(a)$$
 
Donc $\limite{a} f(x)-f(a)=\limite{a} (x-a)f'(a)  =0$.
 


\end{demo}

\section{Théorème des valeurs intermédiaires}

\subsection{Théorème des valeurs intermédiaires}

\begin{thm}[Théorème des valeurs intermédiaires, admis]
 Soit $f$ une fonction continue sur un intervalle $I$ et soit $[a;b]$ un intervalle inclus dans $I$. Alors, pour tout réel $k$ compris entre $f(a)$ et $f(b)$, il existe au moins un réel $c$ de l'intervalle $[a;b]$ tel que $f(x)=k$, c'est  à dire que $f$ prend au moins une fois toutes les valeurs intermédiaires entre $f(a)$ et $f(b)$.
\end{thm}



\begin{rem}
 Avec les hypothèses précédentes, si O est compris entre $f(a)$ et $f(b)$, alors il existe \emph{au moins} un réel $c$ tel que $f(c)=0$. \\

 $f(a)$ n'est pas nécessairement inférieur à $f(b)$.
\end{rem}


\begin{prop}
 Soit $f$ une fonction continue \emph{strictement monotone} sur un intervalle $[a;b]$.Alors, pour tout $k$ compris entre $f(a)$ et $f(b)$, il existe \emph{un unique} réel $c\in[a;b]$ tel que $f(c)=k$.
\end{prop}

\begin{demo}[cas où $f$ est strictement croissante]
 Le théorème des valeurs intermédiaires permet d'affirmer qu'il existe au moins un réel $c$ de $[a;b]$ tel que $f(c)=k$.\par
Si $c'$ est un autre réel de $[a;b]$ tel que $c'<c$, $f$ strictement croissante implique que $f(c')<f(c)$ donc $f(c')\neq k$.\\
De même, si $c'$ est un autre réel de $[a;b]$ tel que $c'>c$, $f$ strictement croissante implique $f(c')>f(c)$ donc $f(c')\neq k$.
\end{demo}

\begin{rem}{Cas particulier de l'équation $f(x)=0$}
 Si $f$ est une fonction continue et strictement monotone sur $[a;b]$ telle que $f(a)\times f(b)<0$, alors l'équation $f(x)=0$ admet une unique solution sur $[a;b]$.
\end{rem}

\subsection{Généralisation à des intervalles quelconques}

\begin{prop}
Le théorème des valeurs intermédiaires admet les généralisations suivantes :

\begin{itemize}
\item[$\bullet$] Si $I=[a;+\infty[$ et $\lim_{x\rightarrow +\infty}f(x)=+\infty$ alors pour tout réel $k\geq f(a)$, l'équation $f(x)=k$ admet une solution dans l'intervalle $[a;+\infty[$.
\item[$\bullet$] Si $I=[a;b[$ et $\lim_{x\rightarrow b}f(x)=l$ alors pour tout réel $k$ compris entre $f(a)$ et $l$ l'équation $f(x)=k$ admet une solution dans $[a;b[$.
\item[$\bullet$] Si $I=[-\infty;b[$, $\lim_{x\rightarrow -\infty}f(x)=l$ et $\lim_{x\rightarrow b}f(x)=-\infty$ alors pour tout réel $k<l$, l'équation $f(x)=k$ admet une solution dans l'intervalle $]-\infty;b[$.
\end{itemize}
\end{prop}

\begin{center}
\begin{tabular}{ccc}
 \begin{minipage}{6cm}
  \begin{tabular}{|c|ccccc|}
   \hline
  $x$ & $a$ & & $c$ & & $+\infty$ \\
   \hline
   &  & &  & & $+\infty$ \\

  & & & & $\nearrow$ &  \\

  $f$ & & & $k$ &  & \\

  & & $\nearrow$ & & & \\

  & $f(a)$ & & & & \\
\hline
  \end{tabular}
 \end{minipage}
&
 \begin{minipage}{6cm}
  \begin{tabular}{|c|ccccc|}
   \hline
  $x$ & $a$ & & $c$ & & $b$ \\
   \hline
   &  & &  &  & $l$ \\

  & & & & $\nearrow$ &   \\

  $f$ & & & $k$ &  & \\

  & & $\nearrow$ & & &  \\

  & $f(a)$ & & & &  \\
\hline
  \end{tabular}
 \end{minipage}
&
 \begin{minipage}{6cm}
  \begin{tabular}{|c|ccccc|}
   \hline
  $x$ & $-\infty$ & & $c$ & & $b$ \\
   \hline
   & $l$ & &  & &  \\

  & & $\searrow$ & & &   \\

  $f$ & & & $k$ &  &  \\

  & &  & & $\searrow$ &  \\

  & & & &   & $-\infty$  \\
\hline
  \end{tabular}
 \end{minipage}
\end{tabular}
\end{center}

\section{Application de la continuité de fonction}
\subsection{Convergence de suite et continuité de fonction}
\begin{thm}[Image d’une suite convergente par une fonction continue] 
Si $f$ est une fonction continue en $l$, limite d’une suite $\suite[u]$, alors la suite $(f(u_n))$ est définie à partir d'un certain rang et  converge vers $f(l)$.

\end{thm}

\begin{demo}
On sait que $f$ est continue en $l$, c'est à dire que, pour toute précision demandé $\epsilon >0$ il existe $\eta_{\epsilon}>0$ tel que pour tout nombre $x$ $\eta_{\epsilon}$-proche de $l$ on a le nombre $f(x)$ qui est $\epsilon$-proche de $f(l)$.\\

D'autre part, $u_n$ converge vers $l$, autrement dit pour n'importe qu'elle précision demandé $\epsilon'>0$ il existe un rang $N_{\epsilon'}$ à partir du quel $u_n$ est $\epsilon'$-proche de $l$.\\

Posons $\epsilon'=\eta_{\epsilon}$, ainsi par définition de $\epsilon'$ on a pour tout $n\se  N_{\epsilon'}$ $u_n$ est $\eta_{\epsilon}$-proche de $l$. Par définition de $\eta_{\epsilon}$ on sait alors que  $f(u_n)$ est $\epsilon$-proche de $f(l)$, et ceux  pour tout $n\se  N_{\epsilon'}$.\\

D'où la convergence de la suite $f(u_n)$ vers $f(l)$.
\end{demo}

\begin{prop}
Soient $I \subset \mathbb{R}$ un intervalle, $f : I \longrightarrow I$ une fonction continue et $\suite[u]$ une suite définie par récurrence par $u_0 \in I $ et $u_{n+1}=f(u_n)$ pour tout $n \in \mathbb{N}*$ et convergeant $l \in I$.\\

Alors $l$ est un point fixe de $f$, c'est à dire que $f(l)=l$.

\end{prop}

\begin{demo}
La convergence de $\suite[u]$ vers $l$ implique deux choses : la première est que la suite $(u_{n+1})_{n\in \mathbb{N}}$ converge aussi vers $l$,
 et la seconde est que $(f(u_n))_{n\in \mathbb{N}}$ converge vers $f(l)$ d'après le théorème précédent, par continuité de $f$ en $l$.
 
  Or $u_{n+1} = f(u_n)$, et on en déduit par unicité de la limite que $l = f(l)$.
\end{demo}



\end{document}
